\documentclass[12pt]{article}

% =========================================================
% Packages
% =========================================================
\usepackage[T1]{fontenc}
\usepackage[utf8]{inputenc}
\usepackage[margin=1in]{geometry}

\usepackage{amsmath}
\usepackage{amssymb}
\usepackage{graphicx}
\usepackage{booktabs}
\usepackage{float}
\usepackage{hyperref}

% =========================================================
% Page settings
% =========================================================
\setlength{\parindent}{0pt}
\setlength{\parskip}{0.7em}

% =========================================================
% Supplemental numbering
% =========================================================
\renewcommand{\thesection}{S\arabic{section}}
\renewcommand{\thesubsection}{S\arabic{section}.\arabic{subsection}}
\renewcommand{\thesubsubsection}
{S\arabic{section}.\arabic{subsection}.\arabic{subsubsection}}

\renewcommand{\thefigure}{S\arabic{figure}}
\renewcommand{\thetable}{S\arabic{table}}
\renewcommand{\theequation}{S\arabic{equation}}

% =========================================================
% Hyperlinks
% =========================================================
\hypersetup{
    colorlinks=true,
    linkcolor=blue,
    urlcolor=blue,
    citecolor=blue
}

% =========================================================
% Document
% =========================================================
\begin{document}

% =========================================================
% Title
% =========================================================

\begin{center}

{\Large\bfseries Supplemental Material}

\vspace{1em}

{\large\bfseries
\textit{MoRemix}: An Online Toolbox for Parameterized Body Motion Manipulation
}

\vspace{1em}

% Yihe Gu

% \vspace{0.3em}

% School of Psychology and Cognitive Science,\\
% East China Normal University

\end{center}

\vspace{1.5em}


% =========================================================
\section{Additional Implementation Details of \textit{MoRemix}}
\label{sup:implementation}
% =========================================================

This section provides additional computational details for the
parent-bone-space representation, motion decomposition and parameterized
manipulation, and motion reconstruction procedures described in the
main manuscript.


% =========================================================
\subsection{Parent-Bone-Space Representation}
\label{sup:parent-space}
% =========================================================

\subsubsection{Skeletal-animation inputs}
\label{sup:skeletal-input}

For skeletal-animation data, editable bones are inferred automatically
from the skeleton tree by retaining bones that possess both a valid local
quaternion track $q_i(t)$ and a non-zero distal child offset.
Let $\mathbf{v}_{i,0}$ denote bone $i$'s rest vector from the proximal
joint (i.e., the joint connected to the parent bone) to the distal joint
in parent-bone coordinates.

Identify the rest vector with the pure quaternion
$\widetilde{\mathbf{v}}_{i,0}=(0,\mathbf{v}_{i,0})$. Given the unit local
quaternion track $q_i(t)$, the implementation rotates this vector directly
by quaternion multiplication. Bone $i$'s parent-bone-space motion sequence is

\begin{equation}
\mathbf{r}_i(t)
=
\operatorname{Vec}\!\left(
q_i(t)\otimes\widetilde{\mathbf{v}}_{i,0}\otimes q_i(t)^{-1}
\right),
\label{eq:sup-skeletal-parent-space}
\end{equation}

where $\otimes$ denotes the Hamilton product and $\operatorname{Vec}(\cdot)$
extracts the vector part of a quaternion. Since $q_i(t)$ is a unit
quaternion, $q_i(t)^{-1}=q_i(t)^{*}$.

The corresponding bone length is

\begin{equation}
L_i
=
\left\|\mathbf{v}_{i,0}\right\|.
\label{eq:sup-skeletal-length}
\end{equation}

For visualization, a point-light-display (PLD) preview can also be
generated by sampling the distal-joint world position at each animation
key time.


% ---------------------------------------------------------
\subsubsection{Point-light-display inputs}
\label{sup:pld-input}
% ---------------------------------------------------------

For PLD data, a CSV file containing marker names (i.e., joint
names), frame indices, and three-dimensional world-space coordinates
$\mathbf{g}_i(t)$ is parsed. Each marker is then assigned a parent marker
$p(i)$ by the user through the binding interface.

According to these parent--child relationships, the toolbox constructs
a layer-wise update order through breadth-first traversal of the parent
graph.

For root markers, where $p(i)=0$, the local and global coordinates are
identical:

\begin{equation}
\mathbf{l}_i(t)
=
\mathbf{g}_i(t).
\label{eq:sup-root-local}
\end{equation}

For non-root markers, the local parent--child vector is

\begin{equation}
\mathbf{l}_i(t)
=
\mathbf{g}_i(t)
-
\mathbf{g}_{p(i)}(t),
\qquad
L_i(t)
=
\left\|\mathbf{l}_i(t)\right\|.
\label{eq:sup-local-vector}
\end{equation}

For markers in the first two update layers, the parent-space basis
coincides with the world-coordinate basis, and therefore

\begin{equation}
\mathbf{r}_i(t)
=
\mathbf{l}_i(t).
\label{eq:sup-first-layers}
\end{equation}

For deeper layers, a moving orthonormal basis

\begin{equation}
\mathbf{B}_i(t)
=
\left[
\hat{\mathbf{x}}_i(t)
\quad
\hat{\mathbf{y}}_i(t)
\quad
\hat{\mathbf{z}}_i(t)
\right]
\label{eq:sup-basis}
\end{equation}

is generated from the orientation of the parent bone. Its first axis is
aligned with the parent local vector:

\begin{equation}
\hat{\mathbf{x}}_i(t)
=
\frac{
\mathbf{l}_{p(i)}(t)
}{
\left\|
\mathbf{l}_{p(i)}(t)
\right\|
}.
\label{eq:sup-x-axis}
\end{equation}

At the first frame,
$\hat{\mathbf{y}}_i(1)$ and $\hat{\mathbf{z}}_i(1)$ are obtained by
orthogonalizing an auxiliary vector against
$\hat{\mathbf{x}}_i(1)$. At subsequent frames, they are updated using
the Rodrigues rotation that aligns the reference
$\hat{\mathbf{x}}$ axis with its current orientation.

The marker trajectory is then expressed in parent-bone-space
coordinates as

\begin{equation}
\mathbf{r}_i(t)
=
\mathbf{B}_i^{\top}(t)
\mathbf{l}_i(t).
\label{eq:sup-pld-parent-space}
\end{equation}

If temporal extension is requested, both skeletal-animation and PLD
clips are repeated by concatenating copies of the corresponding tracks
while removing the duplicated boundary sample between adjacent copies.


% ---------------------------------------------------------
% Optional supplementary figure
% ---------------------------------------------------------

% \begin{figure}[H]
% \centering

% \IfFileExists{Figure-3.1.pdf}{
%     \includegraphics[width=\textwidth]{Figure-3.1.pdf}
% }{
%     \fbox{
%         \parbox[c][5cm][c]{0.9\textwidth}{
%             \centering
%             Placeholder for Figure-3.1.pdf
%         }
%     }
% }

% \caption{
% Detailed computation of parent-bone-space motion
% $\mathbf{r}_i(t)$ for skeletal-animation and PLD inputs.
% For skeletal-animation data, the parent-space bone vector is obtained
% from the local quaternion track and the rest-pose bone vector.
% For PLD data, global joint trajectories are first converted into local
% parent--child vectors and subsequently expressed in a continuously
% updated parent-bone coordinate system.
% }
% \label{fig:sup-parent-space}

% \end{figure}


% =========================================================
\subsection{Motion Decomposition and Parameterized Manipulation}
\label{sup:manipulation}
% =========================================================

Throughout this subsection, $i$ indexes an editable bone in skeletal
animation or a marker--parent segment in PLD data, $t\in\{1,\ldots,T\}$
indexes the $T$ input frames, and $k\in\{1,2,3\}$ indexes a principal
component. Bold symbols denote three-dimensional vectors, an overbar
denotes a temporal mean, a tilde denotes the corresponding zero-mean
component, and a prime denotes an edited quantity.

For each segment $i$, the parent-bone-space trajectory
$\mathbf{r}_i(t)$ over $T$ frames is arranged as a matrix
$\mathbf{X}_i\in\mathbb{R}^{T\times3}$ whose $t$th row is
$\mathbf{r}_i^{\top}(t)$. The toolbox applies uncentered principal
component analysis (PCA), so no temporal mean vector is subtracted from
$\mathbf{X}_i$. Let
$\mathbf{E}_i=[\mathbf{e}_{i,1}\ \mathbf{e}_{i,2}\ \mathbf{e}_{i,3}]$
contain the eigenvectors in descending order of explained variance. The
framewise PCA coordinates are computed explicitly as

\begin{equation}
c_{i,k}(t)
=
\mathbf{e}_{i,k}^{\top}\mathbf{r}_i(t),
\qquad k\in\{1,2,3\}.
\label{eq:sup-pca-score}
\end{equation}

This yields principal axes
$\mathbf{e}_{i,1}$,
$\mathbf{e}_{i,2}$, and
$\mathbf{e}_{i,3}$,
and their corresponding framewise coordinates
$c_{i,1}(t)$,
$c_{i,2}(t)$, and
$c_{i,3}(t)$.

The first two principal axes define the main motion plane (MP), whereas
the third principal axis defines the orthogonal $V_3$ direction.

The \textit{MP Angle} and the corresponding in-plane projection radius
are defined as

\begin{equation}
\theta_i(t)
=
\operatorname{unwrap}
\left[
\operatorname{atan2}
\left(
c_{i,2}(t),
c_{i,1}(t)
\right)
\right],
\label{eq:sup-theta}
\end{equation}

and

\begin{equation}
\rho_i(t)
=
\sqrt{
c_{i,1}^{2}(t)
+
c_{i,2}^{2}(t)
}.
\label{eq:sup-rho}
\end{equation}

Angle unwrapping removes artificial $2\pi$ discontinuities between
adjacent frames.

The MP component is separated into a mean term and a time-varying term:

\begin{equation}
\bar{\theta}_i
=
\frac{1}{T}
\sum_{t=1}^{T}
\theta_i(t),
\qquad
\tilde{\theta}_i(t)
=
\theta_i(t)
-
\bar{\theta}_i.
\label{eq:sup-mp-components}
\end{equation}

Similarly, the $V_3$ component is decomposed as

\begin{equation}
\bar{c}_{i,3}
=
\frac{1}{T}
\sum_{t=1}^{T}
c_{i,3}(t),
\qquad
\tilde{c}_{i,3}(t)
=
c_{i,3}(t)
-
\bar{c}_{i,3}.
\label{eq:sup-v3-components}
\end{equation}


% ---------------------------------------------------------
\subsubsection{Frequency-domain manipulation of the main-plane component}
\label{sup:frequency-manipulation}
% ---------------------------------------------------------

The zero-mean MP component is edited in the frequency domain as

\begin{equation}
\tilde{\theta}_i'(t)
=
\operatorname{Re}\!\left[
\mathcal{F}^{-1}
\left\{
\alpha_i H_i(f)\,
\mathcal{F}\{\tilde{\theta}_i(t)\}
\right\}
\right].
\label{eq:sup-spectrum-edit}
\end{equation}

Here, $\mathcal{F}$ and $\mathcal{F}^{-1}$ denote the FFT and inverse
FFT, respectively; $\operatorname{Re}$ retains the real part of the
result; $f$ denotes frequency; $\alpha_i\geq0$ is the
\textit{MP Amplitude Scale} of segment $i$; and $H_i(f)$ is its filter
gain. The low-pass method retains frequencies below $f_{1,i}$, the
high-pass method retains those above $f_{1,i}$, the band-pass method
retains those in $[f_{1,i},f_{2,i}]$, and the band-stop method attenuates
that interval. Here, $f_{1,i}$ and $f_{2,i}$ are the lower and upper
cutoff frequencies for segment $i$. The stop-band gain is set by the
\textit{Remaining Amplitude} ratio $\eta_i\in[0,1]$, and an optional
transition width $w_i$ replaces an abrupt boundary with a raised-cosine
transition. A sequence of length $T$ is zero-padded to
$N=2^{\lceil\log_2 T\rceil}$, the smallest power of two not less than $T$,
as required by the radix-2 FFT implementation.

The edited mean and time-varying MP components are combined as

\begin{equation}
\theta_i'(t)
=
\tilde{\theta}_i'(t)
+
\bar{\theta}_i+\Delta\theta_i,
\label{eq:sup-edited-theta}
\end{equation}

where $\Delta\theta_i$ is the \textit{MP Mean Addition} applied to
segment $i$.


% ---------------------------------------------------------
\subsubsection{Manipulation of the \texorpdfstring{$V_3$}{V3} component}
\label{sup:v3-manipulation}
% ---------------------------------------------------------

The two $V_3$ controls are applied to the mean and zero-mean components,
respectively. The result is constrained by the parent--child distance:

\begin{align}
c_{i,3}^{*}(t)
&=
\operatorname{clip}
\left(
\gamma_i\tilde{c}_{i,3}(t)+\bar{c}_{i,3}+\beta_i\bar{L}_i,
-L_i(t),L_i(t)
\right),
\label{eq:sup-v3-edit}\\
\rho_i'(t)
&=
\sqrt{
L_i^2(t)
-
\left[c_{i,3}^{*}(t)\right]^2
}.
\label{eq:sup-new-radius}
\end{align}

Here, $c_{i,3}^{*}(t)$ is the constrained edited coordinate on the
third principal axis; $\gamma_i\geq0$ is the \textit{$V_3$ Amplitude
Scale}; $\beta_i$ is the dimensionless \textit{$V_3$ Mean Addition};
$L_i(t)$ is the parent--child distance at frame $t$; and $\bar{L}_i$ is
its temporal mean. For skeletal animation, bone length is fixed, so
$L_i(t)=\bar{L}_i=L_i$. The operator $\operatorname{clip}(x,a,b)$
restricts $x$ to $[a,b]$, and $\rho_i'(t)$ is the resulting in-plane
projection length.

The edited parent-bone-space trajectory is reconstructed in the original
PCA basis:

\begin{equation}
\begin{split}
\mathbf{r}_i'(t)
={}&
\rho_i'(t)
\cos\theta_i'(t)
\mathbf{e}_{i,1}
\\
&+
\rho_i'(t)
\sin\theta_i'(t)
\mathbf{e}_{i,2}
+
c_{i,3}^{*}(t)
\mathbf{e}_{i,3}.
\end{split}
\label{eq:sup-parent-space-reconstruction}
\end{equation}


% ---------------------------------------------------------
\subsubsection{Temporal phase and frequency manipulation}
\label{sup:temporal-manipulation}
% ---------------------------------------------------------

Let $\mathbf{y}_i[m]$, $m\in\{0,\ldots,M_i-1\}$, denote a complete
reconstructed track of segment $i$, where $m$ is the sample index and
$M_i$ is the number of samples after any frequency scaling. The
normalized \textit{Phase Shift} $\phi_i\in[0,1]$ circularly delays the
track by the fraction $\phi_i$ of one sequence:

\begin{equation}
\mathbf{y}_{i,\phi}'[m]
=
\mathbf{y}_i\!\left[
\left(m-\lfloor\phi_iM_i\rfloor\right)\bmod M_i
\right].
\label{eq:sup-phase-direction}
\end{equation}

Thus, $\phi_i=0$ leaves the track unchanged and a positive value shifts
it to later samples without changing its duration. In practice, the
operation is applied to the reconstructed quaternion track for skeletal
animation and to the reconstructed parent-bone-space trajectory before
hierarchical accumulation for PLD data. The floor brackets round the
shift to an integer number of samples, ``$\bmod M_i$'' wraps the sample
index at the sequence boundary, and the subscript $\phi$ identifies the
phase-shifted track.

The \textit{Frequency Scale} $\lambda_i>0$ changes the playback rate by
a time reparameterization:

\begin{equation}
\mathbf{y}_{i,\lambda}'(u)
=\mathcal{I}_i\{\mathbf{y}_i\}(\lambda_i u),
\qquad
0\leq u\leq D_i'=\frac{D_i}{\lambda_i}.
\label{eq:sup-frequency-remapping}
\end{equation}

Here, $u$ is time in the edited track, and $D_i$ and $D_i'$ are the
original and edited durations of segment $i$. The operator
$\mathcal{I}_i$ evaluates the original animation curve at the remapped
time by locating the two neighboring source keyframes and interpolating
between them. Scalar and vector tracks use linear interpolation, whereas
quaternion tracks use normalized spherical linear interpolation. Because
frequency scaling changes the track duration, it also changes the required
number of keyframes. The method estimates the original sampling interval
from the source duration and keyframe count, chooses a new keyframe count
that preserves this interval as closely as possible over the edited
duration, distributes the new keyframes uniformly, and samples the
original curve at their remapped times. The subscript $\lambda$ identifies
the frequency-scaled track. Therefore, $\lambda_i>1$ makes the motion
faster and
$\lambda_i<1$ makes it slower, while every temporal frequency $f$ is
mapped to $\lambda_i f$. If normalized autocorrelation provides a
reliable fundamental-frequency estimate $f_{0,i}$, a requested
\textit{Frequency} $f_i'$ is converted to
$\lambda_i=f_i'/f_{0,i}$; $f_{0,i}$ and $f_i'$ denote the estimated
original and requested fundamental frequencies of segment $i$,
respectively. Otherwise, the user specifies $\lambda_i$ directly.

After every track has been resampled, the common clip duration is set to
the longer of the original clip duration and the longest edited track.
Any shorter track is then extended to this duration by periodic repetition,
rather than by inserting zeros or holding its final pose. This duration
synchronization is performed after, and is distinct from, the resampling
required by the time reparameterization. For PLD data, the same temporal
mapping is propagated from an edited marker to all of its descendants so
that the hierarchy remains synchronized.


% ---------------------------------------------------------
% Optional supplementary figure
% ---------------------------------------------------------

% \begin{figure}[H]
% \centering

% \IfFileExists{Figure-3.2.pdf}{
%     \includegraphics[width=\textwidth]{Figure-3.2.pdf}
% }{
%     \fbox{
%         \parbox[c][5cm][c]{0.9\textwidth}{
%             \centering
%             Placeholder for Figure-3.2.pdf
%         }
%     }
% }

% \caption{
% PCA-based decomposition of parent-bone-space motion.
% The first two principal axes,
% $\mathbf{e}_{i,1}$ and $\mathbf{e}_{i,2}$,
% span the main motion plane, whereas
% $\mathbf{e}_{i,3}$ defines the orthogonal $V_3$ direction.
% }
% \label{fig:sup-decomposition}

% \end{figure}


% =========================================================
\subsection{Motion Reconstruction}
\label{sup:reconstruction}
% =========================================================

\subsubsection{Skeletal-animation reconstruction}
\label{sup:skeletal-reconstruction}

For skeletal-animation data, the edited parent-bone-space direction is

\begin{equation}
\hat{\mathbf{d}}_i'(t)
=
\frac{
\mathbf{r}_i'(t)
}{
\left\|
\mathbf{r}_i'(t)
\right\|
}.
\label{eq:sup-edited-direction}
\end{equation}

Let $\hat{\mathbf{d}}_i(t)$ denote the original parent-bone-space
direction and $q_i(t)$ the original local quaternion.
The toolbox computes the minimal rotation
$q_{\mathrm{dir},i}(t)$ that maps
$\hat{\mathbf{d}}_i(t)$ onto
$\hat{\mathbf{d}}_i'(t)$. Under the active-rotation, column-vector
convention used here, $q_{\mathrm{dir},i}(t)$ is expressed in parent-bone
coordinates and left-multiplies the original local quaternion:

\begin{equation}
q_i'(t)
=
q_{\mathrm{dir},i}(t)
q_i(t).
\label{eq:sup-quaternion-update}
\end{equation}

This preserves the original axial twist while enforcing the edited bone
direction. The resulting quaternion sequence then overwrites the
corresponding animation track in the current clip.


% ---------------------------------------------------------
\subsubsection{Point-light-display reconstruction}
\label{sup:pld-reconstruction}
% ---------------------------------------------------------

For PLD data, reconstruction proceeds hierarchically according to the
previously defined update order.

For markers in the first two update layers, the edited local trajectory
is set directly as

\begin{equation}
\mathbf{l}_i'(t)
=
\mathbf{r}_i'(t).
\label{eq:sup-pld-first-layers-reconstruction}
\end{equation}

For deeper layers, a new moving basis
$\mathbf{B}_i'(t)$ is recomputed from the edited parent local trajectory
using the same Rodrigues-update procedure. The edited local vector is

\begin{equation}
\mathbf{l}_i'(t)
=
\mathbf{B}_i'(t)
\mathbf{r}_i'(t).
\label{eq:sup-pld-local-reconstruction}
\end{equation}

Global marker trajectories are subsequently reconstructed recursively:

\begin{equation}
\mathbf{g}_i'(t)
=
\begin{cases}
\mathbf{l}_i'(t),
& p(i)=0, \\[6pt]
\mathbf{g}_{p(i)}'(t)
+
\mathbf{l}_i'(t),
& p(i)\neq0.
\end{cases}
\label{eq:sup-global-reconstruction}
\end{equation}

% =========================================================
\section{Additional Information for the Validation and Demonstration Experiment}
\label{sup:demonstration}
% =========================================================

\subsection{Action Categories and Analyzed Bones in the Validation Experiment}
\label{sup:validation-actions}

The validation analysis included six everyday action categories, grouped into
full-body, upper-limb, and lower-limb actions. For each category, the analysis
focused on the major bones primarily involved in the corresponding action,
rather than applying an identical set of bones to all actions. The action
categories and analyzed bone units are summarized in
Table~\ref{tab:sup-validation-bones} using original bone names in AMASS datasets (Figure.~S1).

% Table 1
\begin{table}[H]
\centering

\caption{
Mapping between body-part categories and bones in the SMPL-based skeleton. 
}
\label{tab:sup-validation-bones}

\begin{tabular}{lll}
\toprule
Body part &  Bones \\
\midrule
Body  & Spine1, Spine2, Spine3  \\
Head  & Neck  \\
Upper arm & L/R\_Shoulder \\
Lower arm & L/R\_Elbow \\
Upper leg & L/R\_Hip \\
Lower leg & L/R\_Knee \\
\bottomrule
\end{tabular}

\end{table}

% Table 2
\begin{table}[H]
\centering

\caption{
Action categories and major bone units analyzed in the validation experiment.
}
\label{tab:sup-validation-bones}

\begin{tabular}{lll}
\toprule
Action group & Action category & Body-part (L\&R)\\
\midrule
Full-body  & walk      & Body, Head, Upper arm, Lower arm, Upper leg, Lower leg \\
Full-body  & run       & Body, Head, Upper arm, Lower arm, Upper leg, Lower leg \\
Upper-limb & wave      & Upper arm, Lower arm \\
Upper-limb & stretch   & Upper arm, Lower arm \\
Lower-limb & side step left & Upper leg, Lower leg \\
Lower-limb & kick      & Upper leg, Lower leg \\
\bottomrule
\end{tabular}

\end{table}

\begin{figure}[H]
    \centering
    \includegraphics[width=0.9\textwidth]{SMPL_bonelist.png}
    \caption{The bone name list of SMPL-based skeleton used in AMASS datasets}
    \label{fig:sup-figure1}
\end{figure}


\subsection{Body-Part Definitions in the Demonstration Experiment}
\label{sup:demonstration-bodyparts}

In the demonstration experiment, the six manipulated body-part categories were
defined at the bone level. Bilateral limb categories included the corresponding
left and right bones, whereas the spine and head categories each referred to the
corresponding central bone. The action
categories and manipulated bone units are summarized in
Table~\ref{tab:sup-demonstration-bones} using original bone names in MetaHuman Skeleton (Figure.~S2).

\begin{table}[H]
\centering

\caption{
Mapping between body-part categories and bones in the MetaHuman skeleton.
}
\label{tab:sup-demonstration-bones}

\begin{tabular}{ll}
\toprule
Body-part category & Included bones \\
\midrule
Body     & pelvis, spine\_01-05 \\
Head      & neck\_01-02 \\
Upper arm & upper arm\_l/r \\
Lower arm & lower arm\_l/r \\
Upper leg & upper leg\_l/r \\
Lower leg & lower leg\_l/r \\
\bottomrule
\end{tabular}

\end{table}

\begin{figure}[H]
    \centering
    \includegraphics[width=0.9\textwidth]{MetaHuman_Bonelist.png}
    \caption{The bone name list of MetaHuman Skeleton}
    \label{fig:sup-figure1}
\end{figure}

Because of rendering limitations of the online toolbox, the provided stimulus animation file includes only the skeletal structure and animation tracks, without an associated character mesh. For the demonstration experiment, the same skeletal animation was retargeted to six different Asian MetaHuman character meshes (Figure.~S3), which were used to render the experimental stimuli.

\begin{figure}[H]
    \centering
    \includegraphics[width=0.9\textwidth]{MetaHumanMesh.png}
    \caption{Asian MetaHuman characters used in the demonstration experiment}
    \label{fig:sup-figure1}
\end{figure}

\subsection{Psychometric Fitting and Missing-Value Summary}
\label{sup:psychometric-fitting}

Individual response data were fitted using logistic psychometric
functions implemented in the Palamedes toolbox for MATLAB, as described
in the main manuscript. Separate psychometric functions were fitted for
the lower and upper amplitude-manipulation ranges of each body part.

Threshold estimates were defined as the manipulation levels
corresponding to a 50\% probability of judging the stimulus as natural.
Fits that did not yield a valid threshold estimate were coded as missing
and were excluded from the corresponding threshold summary.

Table~\ref{tab:sup-missing-values} reports the number of missing
threshold estimates for each bone and manipulation direction.

\begin{table}[H]
\centering

\caption{
Number of missing threshold estimates in the demonstration experiment (full sample size: 30).
}
\label{tab:sup-missing-values}

\begin{tabular}{lcc}
\toprule
Bone &
Lower threshold &
Upper threshold \\
\midrule
Spine     & 4 & 0 \\
Head      & 0 & 12 \\
Upper arm & 6 & 4 \\
Lower arm & 3 & 10 \\
Upper leg & 0 & 0 \\
Lower leg & 1 & 1 \\
\bottomrule
\end{tabular}

\end{table}



% =========================================================
% End
% =========================================================

\end{document}
